\begin{theorem}\label{night:P1.19}
\StatementRecord{night:P1.19}{proved}{}{\cref{night:P1.11}, \cref{night:P1.12}, \cref{night:P1.9}; \cite[Theorem~3.43, p.~254]{Hatcher2002}}
There is an independent computation from the full global cellular
complex.  Its chain ranks in degrees zero through eight are
\[
 (4,26,87,162,200,156,79,23,3),
\]
and its differentials, after the explicit integral chain equivalences
below, are the matrices in
\path{code/current_audits/audit_corrections/code/reconstruct_V_global_06.json}.
The accompanying program reconstructs them from the displayed marked
maps; the separate saved-output verifier supplies their integral
cycle and incoming-map checks and the basis comparisons below.
Their integral Smith calculation gives
\[
 H_k(M;\mathbb Z)=
 \begin{cases}
 \mathbb Z,&k=0,8,\\
 \mathbb Z^4,&k=4,\\
 0,&\text{otherwise}.
 \end{cases}
\]
This calculation uses local Poincar\'e--Lefschetz duality of the
smooth cusp cap to determine its secondary attaching class.  It does
not use global Poincar\'e duality or the homology conclusion of P1.12.
In particular, the map obtained by setting the whole cusp suspension
chain term to zero is not the actual relative cusp model; its precise
missing term is accounted for here.
\end{theorem}
\begin{proof}
\begin{enumerate}[label=\arabic*.]
\item We retain the actual complement complex and finite-cap maps of P1.12.
The issue left by that argument was the full cusp boundary map.
We first determine its homology maps, then its secondary chain term,
and finally give the explicit integral reduction used by the program.

\item Write the cusp lattice as $A\oplus B$, with ordered generators
$a_1,a_2,a_3,b_1,b_2,b_3$ and $Nb_i=a_i$.
The boundary is the mapping torus of
\[
 (b,a)\longmapsto(b,a+b)
\]
on $T_B^3\times T_A^3$.
For the pure toric model, its retraction to the central moment model
has the following explicit description.  On the dual cell of a Kuhn
simplex $\sigma$, choose any vertex $v$ of $\sigma$ and put
\[
 (b,a,\theta)\longmapsto
   [b,a+\theta(b-v)],\qquad0\leq\theta\leq1.
\tag{1.19.1}
\]
The brackets denote the angular quotient by $L_\sigma$.
Changing $v$ changes the displayed angular coordinate by a real
multiple of a vector in $L_\sigma$, so the map is unchanged.
Replacing $b$ by $b+n$ also replaces $v$ by $v+n$; hence the formula
is periodic.  At $\theta=1$ the integer translation $v$ disappears
on the angular torus, and the map agrees with the cusp monodromy
identification.  These observations prove that (1.19.1) is well-defined
on the boundary mapping torus.  In a monomial chart of P1.9 it is
exactly the operation of forgetting the arguments of the vanishing
coordinates, with the argument of $q$ recorded by $\theta$.
Thus it is the actual boundary retraction.  The deformation removing
the bounded period units in P1.11 transports this map by a homotopy
and preserves its homology class.

\item We prove that the suspension classes have zero image in cusp homology.
Let $D$ be the derivation of $\bigwedge^\bullet(A\oplus B)$ extending
$Nb_i=a_i$, $Na_i=0$.  Then $\bigwedge M_0=\exp D$ and
$\ker(\bigwedge M_0-I)=\ker D$ over $\mathbb Q$.
An integral vector in this kernel has the same property.
The summand with $p$ base generators and $q$ angular generators is
sent by $D$ to the summand with $(p-1,q+1)$; consequently each
homogeneous summand of an invariant vector is separately killed.
Let $E$ replace $a_i$ by $b_i$. In the orthonormal exterior basis
$E=D^t$ and $[D,E]=(q-p)I$ on that summand.  For $Dv=0$,
\[
 (q-p)\|v\|^2
 =\langle[D,E]v,v\rangle=\|Ev\|^2\geq0.
\]
Thus an invariant $k$-cycle has base degree at most $\lfloor k/2\rfloor$.

\item This assertion about cycles can be realized geometrically, rather
than only in a rational associated graded group.  Give $T_B^3$ its
minimal product CW structure and put
$F_p=T_A^3\times(T_B^3)^{(p)}$.
Its cellular differential is zero, so
$H_*(F_p;\mathbb Z)\to H_*(T^6;\mathbb Z)$ is injective.
The shear preserves $F_p$ and admits cellular approximation through
maps preserving this filtration.  The corresponding mapping-cylinder
cells give the integral complex
$\operatorname{Cone}(\bigwedge M_0-I)$.
An invariant cellular vector of base degree at most $p$ therefore
has a suspension lift whose entire projection to $T_B^3$ has dimension
at most $p$.  Equivalently, its closing boundary already bounds in
$F_p$, by the just-proved injectivity.  This construction fixes the
suspension representatives used in the Wang complex and is compatible
with the based product-cell comparison in P1.12.

\item Formula (1.19.1) preserves $b$, so the image of such a lift also has
base projection of dimension at most $p$.  It represents filtration
at most $p$ in the dual-cell filtration of P1.11.  Here is a direct
cellular justification of this last step.  In every dual cell of
dimension greater than $p$, choose an interior point outside that
projection.  Radial retraction from this point pushes the projected
cycle to the cell boundary.  On the central moment model it lifts
by keeping the angular coordinate fixed; on the boundary the angular
quotient is the canonical further quotient.  The lift is continuous
and fixes the boundary.  Performing these retractions in descending
cell dimension gives a homologous cycle on the dual $p$-skeleton.
There are finitely many cells modulo periods, so this operation
terminates.

\item The integral $E_2=E_\infty$ table of P1.11 has nonzero entries only
at angular degree $q\leq d$, where $d$ is the base-cell degree.
For a suspension of a $k$-cycle, the total degree is $k+1$, whereas
$d\leq p\leq\lfloor k/2\rfloor$. Its angular degree would therefore
be $k+1-d>d$, where that table is zero.  The suspension image
vanishes.  On the fibre summand the map is precisely the marked
collapse $\kappa_k$ reconstructed in
\path{code/current_audits/audit_corrections/code/reconstruct_V_cusp_06.json}.
Its source basis is $(a_1,a_2,a_3,b_1,b_2,b_3)$; the actual matrix $S$
from P1.7 changes it to the original lattice by
$\kappa_k^{\rm orig}=\kappa_k\bigwedge^kS^{-1}$.
These transported matrices are included in the global reconstruction.  Thus the homology maps agree
with those of the chain map $f_0=(\kappa,0)$, in every degree.

\item Agreement on homology does not identify these two chain maps.
We now determine the remaining difference.  Integral Smith reduction
splits a finite free complex into free homology generators, contractible
unit pairs, and two-term pairs with differential equal to a torsion
invariant factor.  The cusp homology is free.  The boundary homology
has just one torsion group, $\mathbb Z/2$ in degree three.
Consequently, after maps on homology have been fixed, the difference
of two maps from the boundary complex to the cusp complex is
represented by one vector
\[
 \varepsilon\in H_4(C;\mathbb Z)/2H_4(C;\mathbb Z)
                 \cong(\mathbb Z/2)^{10}.
\tag{1.19.2}
\]
This follows directly by applying a map to the upper generator of
the two-term pair $\mathbb Z\xrightarrow{2}\mathbb Z$; changing a
chain homotopy changes its image by twice an integral vector.
There are no other secondary terms.  Thus no rational formality or
unproved equivariant formality is used at this point.

\item In the marked cusp degree-four basis, the primitive functional
\[
 \psi=(1,-1,1,1,-1,-1,1,0,0,0)
\tag{1.19.3}
\]
annihilates $\kappa_4$ and identifies its cokernel with $\mathbb Z$.
The exact multiplication $\psi\kappa_4=0$ and the primitive Smith
factors are checked in the global reconstruction.  Setting
$\varepsilon=0$ gives relative cusp homology
$H_4(C,\partial)=\mathbb Z^{10}\oplus\mathbb Z/2$.
The two-term Smith-pair splitting in the preceding item gives this
split extension directly.  The saved-output verifier also constructs
the relative cone and checks its degree-four integral homology.
More intrinsically, the relative exact sequence in this degree is
an extension of $\mathbb Z^9\oplus\mathbb Z/2$ by the cokernel
$\mathbb Z$ detected by $\psi$.  Its torsion extension class is
exactly $\psi\varepsilon$ modulo two.

\item The actual cusp neighborhood is a compact oriented smooth
8-manifold with boundary: take the inverse image of a sufficiently
small closed disc, whose boundary circle avoids the critical value.
Let $\rho$ be its retraction onto $C$.  Pullback by $\rho$, followed
by cap product with the relative fundamental class of this actual
oriented cusp cap, is the isomorphism of
\cite[Theorem~3.43, p.~254]{Hatcher2002}, with $n=8$, $k=4$,
$R=\mathbb Z$, $A=\varnothing$ and $B=\partial$.
The already computed free cusp homology therefore gives
\[
 H_4(\text{cusp cap},\partial;\mathbb Z)
       \cong H^4(C;\mathbb Z)=\mathbb Z^{10}.
\]
The preceding extension must therefore be nonsplit on its torsion
summand.  This determines the required secondary class for the global
calculation:
\[
                 \psi\varepsilon=1\pmod2.
\tag{1.19.4}
\]
This also proves the stated failure of the wholly zero suspension
chain map as a relative cusp model.

\item We next give the integral equivalence which reduces every class
satisfying (1.19.4) to one displayed matrix.  Write the degree-four
homology gluing map in the plus-sign convention as
\[
 j_4:H_4(\partial)\longrightarrow H_4(Y)\oplus H_4(C).
\]
Multiplication by $-1$ on the cusp chain summand is the chain
isomorphism from the minus-sign convention in the exported whole
diagram to this convention.  The matrix $j_4$ is $22$ by $18$,
has rank eighteen, and all its nonzero Smith factors are one.
Its cokernel is free of rank four.  The saved-output verifier exports
the actual boundary cycle basis and complement homology projection
used to form $j_4$.  Write $U_4j_4V_4$ for its exported Smith reduction;
take the last four rows of $U_4$ and interchange their first and
fourth rows.  This is the explicitly verified unimodular cokernel
adapter.  In the resulting cokernel basis the map from $H_4(C)$ is
\[
 \begin{pmatrix}0\\0\\0\\\psi\end{pmatrix}.
\tag{1.19.5}
\]
This is a direct integer matrix calculation, independent of the
homology of $M$.

\item Let $e_4$ denote the fourth coordinate vector of $\mathbb Z^{10}$,
so $\psi e_4=1$.  The nine columns
$e_i-\psi_i e_4$, $i\ne4$, form an integral basis $K$ of $\ker\psi$;
$[e_4,K]$ is unimodular.  The saved-output verifier gives the exact
unimodular change from the Smith kernel basis to these nine columns,
and exports the resulting integral lifts $W$ satisfying
\[
 j_4W=\begin{pmatrix}0\\K\end{pmatrix}.
\]
It also gives actual boundary cycles $X$ representing these lifts,
and integral complement chains $T$ with
\[
 d_YT=j_YX,\qquad f_0X=K.
\tag{1.19.6}
\]
Every entry is checked by multiplication in the program.

\item For clarity, the chain automorphisms implementing the reduction
are equally explicit.  Let $r$ be the exported integral degree-four
cocycle representing the torsion generator in $H^4(\partial;\mathbb Z)$.
In the ordered boundary group
$B_{0,4}=\bigwedge^4\Lambda\oplus\bigwedge^3\Lambda$, with each
exterior basis lexicographically ordered in the original lattice,
it is
\[
r=e_{29}^{t}+e_{30}^{t}+e_{31}^{t}
  -e_{32}^{t}-e_{33}^{t}-e_{34}^{t}+e_{35}^{t}.
\]
The verifier checks that it is not an integral coboundary by the
augmented rank obstruction modulo two.  It satisfies
\[
 r d_{\partial,5}=0,\qquad rX=0,
 \qquad 2r=t\,d_{\partial,4},
\tag{1.19.7}
\]
where the integral row $t$ is exported as well.
For a column $v\in\mathbb Z^9$, the endomorphism
$\eta=(Xv)r$ of degree-four boundary chains, zero in all other
degrees, is a chain map and has $\eta^2=0$.
Thus $I+\eta$ is an integral chain automorphism with inverse
$I-\eta$. Precomposing the attaching map by it changes the
secondary vector by $Kv$, up to the explicitly given homotopy
$(Tv)r$ in the complement, by (1.19.6).
The last equation in (1.19.7) removes an even multiple of $e_4$
by the homotopy $e_4t$.  Decomposing $\varepsilon$ in the
unimodular basis $[e_4,K]$, equation (1.19.4) therefore reduces
the actual global attaching map to the canonical correction
\[
 f_4=f_{0,4}+e_4r,
 \qquad f_k=f_{0,k}\quad(k\ne4).
\tag{1.19.8}
\]
These formulas prove the chain equivalence; they do not merely
compare the ranks of the two complexes.

\item The reconstruction program inserts (1.19.8) into the complete
main-piece/three-cap mapping cone, retaining every finite-cap
differential and every boundary map of P1.12.  It checks $d^2=0$
over $\mathbb Z$ and exports every differential.  The corrected case
is the actual global model by the preceding explicit chain
equivalence.

The separate saved-output verifier checks the integral inverse
matrices and cycle bases, the incoming matrices in those bases, and
their Smith left and right matrices.  It verifies the outgoing ranks
by their displayed integral kernels and exact rank witnesses modulo
$1009$.  It also exports the actual $W,X,T$ maps in the literal $K$
basis, checks every identity in (1.19.6)--(1.19.7), and verifies both
relative cusp degree-four groups.  The even case is recomputed as
an unused diagnostic and gives the same global homology; its local
relative torsion retains the distinction needed in item~9.
The differential ranks in degrees one through eight are
\[
                 (3,23,64,98,98,58,21,2).
\]
All nonzero incoming Smith factors in the cycle lattices are one.
Subtracting incoming and outgoing ranks from the displayed chain
ranks gives $(1,0,0,0,4,0,0,0,1)$, with no torsion.
The exact saved files are
\path{code/current_audits/audit_corrections/code/reconstruct_V_global_06.py},
\path{code/current_audits/audit_corrections/code/reconstruct_V_global_06.json},
\path{code/current_audits/audit_corrections/code/verify_V_global_saved_07.py}
and its JSON output.  This completes the full cellular comparison
without using the global duality conclusion of P1.12.

\end{enumerate}
\end{proof}
